
\documentclass{article}
\usepackage{colortbl}
\usepackage{makecell}
\usepackage{multirow}
\usepackage{supertabular}

\begin{document}

\newcounter{utterance}

\centering \large Interaction Transcript for game `cladder', experiment `full\_v1.5\_default', episode 4893 with Qwen3.6{-}27B.
\vspace{24pt}

{ \footnotesize  \setcounter{utterance}{1}
\setlength{\tabcolsep}{0pt}
\begin{supertabular}{c@{$\;$}|p{.15\linewidth}@{}p{.15\linewidth}p{.15\linewidth}p{.15\linewidth}p{.15\linewidth}p{.15\linewidth}}
    \# & \multicolumn{2}{c}{Player} && \multicolumn{2}{c}{Game Master} \\
    \hline

    \theutterance \stepcounter{utterance}  
    & & & \multicolumn{4}{p{0.6\linewidth}}{
        \cellcolor[rgb]{0.9,0.9,0.9}{
            \makecell[{{p{\linewidth}}}]{
                \texttt{\tiny{[P1$\langle$GM]}}
                \texttt{You are an expert in causal inference. The following question is not a typical commonsense query, but rather a meticulously designed question created by a professor specializing in causal inference, intended to assess the students' mastery of the course content.} \\
\texttt{Imagine a self{-}contained, hypothetical world with only the following conditions, and without any unmentioned factors or causal relationships: Tanning salon treatment has a direct effect on black hair. Going to the beach has a direct effect on black hair. The overall probability of tanning salon treatment is 99\%. For people not using tanning salon treatments, the probability of black hair is 66\%. For people who went to tanning salons, the probability of black hair is 24\%. Is black hair less likely than blond hair overall?} \\
\texttt{Start your answer with 'Yes' or 'No', followed by additional reasoning or evidence to support your explanation.} \\
            }
        }
    }
     \\ \\

    \theutterance \stepcounter{utterance}  
    & \multicolumn{4}{p{0.6\linewidth}}{
        \cellcolor[rgb]{0.561,0.737,0.561}{
            \makecell[{{p{\linewidth}}}]{
                \texttt{\tiny{[P1$\rangle$GM]}}
                \texttt{No} \\
\\ 
\texttt{To determine whether black hair is less likely than blond hair overall, we must calculate the total probability of having black hair, denoted as \$P(B)\$.} \\
\\ 
\texttt{**1. Define the Variables and Given Probabilities:**} \\
\texttt{*   Let \$T\$ be the event of receiving tanning salon treatment.} \\
\texttt{*   Let \$B\$ be the event of having black hair.} \\
\texttt{*   Let \$B\^{}c\$ be the event of not having black hair (which, in this binary context of hair color likelihoods implied by the question, corresponds to blond hair or non{-}black hair).} \\
\\ 
\texttt{From the problem description:} \\
\texttt{*   \$P(T) = 0.99\$} \\
\texttt{*   Therefore, \$P(T\^{}c) = 1 {-} 0.99 = 0.01\$ (probability of not using tanning salon treatments).} \\
\texttt{*   \$P(B | T\^{}c) = 0.66\$ (probability of black hair given no tanning salon treatment).} \\
\texttt{*   \$P(B | T) = 0.24\$ (probability of black hair given tanning salon treatment).} \\
\\ 
\texttt{*Note: The statement "Going to the beach has a direct effect on black hair" is mentioned as a condition in the world, but no specific probabilities or conditional dependencies regarding the beach are provided. In causal inference problems of this type, if specific numerical data for a variable is not given, and the question asks for} \\
            }
        }
    }
    & & \\ \\

    \theutterance \stepcounter{utterance}  
    & & & \multicolumn{2}{p{0.3\linewidth}}{
        \cellcolor[rgb]{0.9,0.9,0.9}{
            \makecell[{{p{\linewidth}}}]{
                \texttt{\tiny{[GM$|$GM]}}
                \texttt{No} \\
            }
        }
    }
    & & \\ \\

    \theutterance \stepcounter{utterance}  
    & & & \multicolumn{2}{p{0.3\linewidth}}{
        \cellcolor[rgb]{0.9,0.9,0.9}{
            \makecell[{{p{\linewidth}}}]{
                \texttt{\tiny{[GM$|$GM]}}
                \texttt{yes} \\
            }
        }
    }
    & & \\ \\

    \theutterance \stepcounter{utterance}  
    & & & \multicolumn{2}{p{0.3\linewidth}}{
        \cellcolor[rgb]{0.9,0.9,0.9}{
            \makecell[{{p{\linewidth}}}]{
                \texttt{\tiny{[GM$|$GM]}}
                \texttt{game\_result = LOSE} \\
            }
        }
    }
    & & \\ \\

\end{supertabular}
}

\end{document}
